
We evaluate the Rust prototype (Section~\ref{sec:rust-runtime}) on three questions:
does replay reproduce a recorded run (E1), does recovery from a crash at any point preserve
what was journaled and complete the run (E2), and what does journaling cost (E3)?
These experiments test that the implementation follows the record, replay, and resume
semantics for a single agent, matching the scope of the formal results, and they use a seeded
stand-in for the LLM.

\paragraph{Setup.}
The agent is a small loop with two branches.
It draws \texttt{classify} from $\{\texttt{tool}: 5/6,\ \texttt{answer}: 1/6\}$.
On \texttt{tool} it draws \texttt{tool-arg} from $\{\texttt{arg-A}: 3/5,\ \texttt{arg-B}: 2/5\}$
and performs a durable effect \texttt{commit(arg)} under the three-phase write-ahead
discipline; on \texttt{answer} it draws \texttt{answer}, whose distribution has a single
outcome.
A recorded run therefore has two or four journal entries.
The oracle is \texttt{SeededOracle}, a categorical sampler driven by SplitMix64, and the journal
is \texttt{InMemoryJournal}.
All seeds come from a fixed SplitMix64 stream, so every experiment is reproducible.
Measurements were taken on an Apple M5 Pro with 48\,GB of memory, macOS~27.0.1, and
rustc~1.93.0, in release builds.
Each experiment is an integration test in the accompanying repository
(\texttt{tests/eval\_e1.rs}, \texttt{eval\_e2.rs}, \texttt{eval\_e3.rs}) run with \texttt{cargo test --release --test <name> -- --nocapture}.

\subsection{E1: Replay Determinism}
\label{sec:eval-e1}

For each of $N = 1{,}000$ trials we record a run with one seed, then replay the complete
journal with a fresh oracle seeded independently.
We compare the traces (each label carries its journal key and the realized value) and the
final outputs.
Theorem~\ref{thm:replay}(i)--(ii) predicts identical traces in every trial.

\paragraph{Result.}
All $1{,}000$ replays reproduced the recorded trace and output exactly.
Every trial made at least one oracle draw, $810$ took the branch with the durable effect, and
the replay seed differed from the record seed in every trial, so agreement comes from the
journal rather than from the oracle.
Because replay never consults the live oracle, this outcome is expected by construction; the
experiment checks that the implementation conforms to \textsc{[O-rep]} and \textsc{[Eff-rep]}.

\subsection{E2: Crash Recovery at Every Journal Position}
\label{sec:eval-e2}

For each of $1{,}000$ recorded runs, and for every crash point $k$ from $0$ to the journal
length, we keep only the first $k$ journal entries, then resume in replay mode with a fresh
oracle whose seed differs from the recording seed.
If the crash point separates an effect's intent from its committed record, the resumed run
stops at the dangling intent; we then apply Class-I recovery (re-perform with the journaled
argument and commit) and resume again.
For each resume we check that
(a)~the run completes with an output and leaves no dangling intent;
(b)~the first $k$ journal entries are unchanged and every trace label whose entry survived is
reproduced;
(c)~every oracle value lies in the support of its distribution and journal cursors stay dense;
and (d)~an effect whose committed record survived is not performed again.

\paragraph{Result.}
The sweep made $4{,}664$ resumes.
All $4{,}664$ completed and reproduced their surviving prefix, $832$ of them crashed inside an
effect's write-ahead window and were recovered, and no surviving committed effect was performed
again.

\paragraph{Conditional law.}
Corollary~\ref{cor:resume}(ii) says that the resumed continuation follows the original law conditioned on the
replayed prefix.
We check this on the prefix \texttt{classify = tool}: across $20{,}000$ resumes with fresh seeds,
the re-sampled \texttt{tool-arg} was \texttt{arg-A} with frequency $0.5989$, against
$3/5 = 0.6$ (standard error $\approx 0.0035$).

\paragraph{Compensation.}
After a crash inside the write-ahead window followed by a compensating recovery (Class~C),
replaying the journal twice stops both times at the compensated effect, produces no output,
performs no effect, and appends nothing, as the write-ahead design intends.

\paragraph{A defect found by E2.}
An earlier version of the driver stayed in replay mode after \textsc{[O-resume]} fell through to
recording, so the next step re-read the entry just appended as if it were history and could
produce out-of-support values.
The original test suite missed this because its only crash test took the branch whose
distribution has a single outcome; the sweep above fails on that version and passes on the
corrected one.

\subsection{E3: Journaling Overhead}
\label{sec:eval-e3}

\paragraph{Append latency.}
We append $n \in \{10^3, 10^4, 10^5, 10^6\}$ oracle-draw events to an in-memory journal, after
one warm-up pass, and report the mean time per append.
Over six runs, every measurement fell between $88$ and $145$\,ns per append, with no growth in
$n$, consistent with amortized constant-time append; the journal holds exactly one entry per
step.
These are single timed passes per size, not a statistical benchmark.

\paragraph{Replay versus re-execution.}
To model a networked LLM call we inject a latency of $1$\,ms into every oracle draw.
A recorded run then takes about $2.51$\,ms, while replaying it takes $0.5$--$0.8\,\mu$s, a
speedup of about $3{,}200\times$ to $4{,}900\times$ across six runs.
The speedup is dominated by the injected latency: replay performs no oracle calls, so the
time saved per replayed draw is the latency of the call itself.

\paragraph{Not measured.}
We did not measure the cost of journaling relative to a journal-free execution mode, which the
prototype lacks, or the cost of persisting the journal (for example, \texttt{fsync} on every
append), since no persistent backend exists yet.

\subsection{Threats to Validity}
\label{sec:eval-threats}

\paragraph{Internal.}
The seeded oracle stands in for an LLM.
The experiments show that the implementation follows the record/replay semantics, not how real
models behave.
Crashes are simulated by truncating an in-memory journal; torn writes, \texttt{fsync} behavior,
and storage failures are not exercised.

\paragraph{Construct.}
Trace comparison in E1 and E2 uses labels that carry cursors, prompts, and responses, as in
Theorem~\ref{thm:replay}.
E2 checks the conditional law on one prefix and one prompt, not in general.

\paragraph{External.}
All measurements come from one machine and a single in-memory backend, and the workload is a
small agent; results for persistent backends, larger agents, or many concurrent agents may
differ.

\begin{table}[t]
\centering
\caption{Summary of the evaluation.}
\label{tab:eval-summary}
\fitwidth{%
\begin{tabular}{@{}lllp{5.2cm}@{}}
\toprule
\textbf{ID} & \textbf{Question} & \textbf{Result relates to} & \textbf{Outcome} \\
\midrule
E1 & Replay determinism & Theorem~\ref{thm:replay} & $1{,}000/1{,}000$ identical replays \\
E2 & Crash at every position & Corollary~\ref{cor:resume}; side effects & $4{,}664/4{,}664$ resumes correct; no duplicate effects; conditional law matches \\
E3 & Journaling overhead & implementation & 88--145\,ns per append, flat in $n$ \\
\bottomrule
\end{tabular}%
}
\end{table}
